% SOURCE_AUTHORITY: sga5_fr_workpass.tex, lines 453--469 (LF-normalized unit).
Abordemos ahora las ``\emph{fórmulas de intercambio}''.

No volveremos sobre la noción de dimensión Tor de un complejo (SGA~A~XVII~4.1.9)~(*)\footnote{(*) Véase también (SGA~6~I~5).}. Retengamos solamente que, si $X$ es un esquema y $G \in \ob \D^{-}(X)$, la relación $\underline{\Tor}_i(F,G) = 0$ para $i \geqslant N$ y todo $A_X$-módulo $F$ equivale a la misma relación con $F$ constructible (gracias a SGA~A~IX~2.9 y a la conmutatividad de $\mathop{\otimes}\limits^{\mathrm{L}}$ con los límites inductivos); por tanto, no es necesario introducir una noción de ``cuasi-dimensión Tor''.

\subsection*{Lema 1.9 (cf. [H] II 4.3)}
\begin{enumerate}
\item[a)] Sean $X$ un esquema, $F \in \ob \D_c(X)$ y $G \in \ob \D_c^{-}(X)$. Se tiene $F \mathop{\otimes}\limits^{\mathrm{L}} G \in \ob \D_c(X)$ si $F \in \ob \D_c^{-}(X)$ o si $G$ tiene dimensión Tor finita.
\item[b)] Sean $f : X \to Y$ un morfismo de esquemas y $G \in \ob \D_c(Y)$. Entonces $f^{*}(G) \in \ob \D_c(X)$.
\end{enumerate}

\emph{Demostración.} La afirmación b) se deduce inmediatamente de (SGA~A~IX~2.4). Demostremos a). Cualquiera de las dos hipótesis sobre $(F,G)$ implica que la sucesión espectral
\[
\E_2^{pq} = \sum_{q_1+q_2=q} \underline{\Tor}_{-p}(\H^{q_1}(F), \H^{q_2}(G)) \Rightarrow \H^{*}(F \mathop{\otimes}\limits^{\mathrm{L}} G)
\]
es birregular. Por tanto, basta demostrar a) cuando $F$ y $G$ están concentrados en grado $0$, es decir, son $A_X$-módulos constructibles. Como la cuestión es local, podemos suponer que $X$ es noetheriano. Entonces $F$ admite (SGA~A~IX~2.7) una resolución por la izquierda $F'$, cuyos términos $F'^i$ son de la forma $A_{U,X}$, con $U$ étale y de tipo finito sobre $X$. Se tiene $F \mathop{\otimes}\limits^{\mathrm{L}} G \xleftarrow{\sim} F' \otimes G$, y la conclusión se sigue de que los $A_{U,X} \otimes G$ son constructibles.
